\documentclass{erciyesfk}
\usepackage[utf8]{inputenc}
\usepackage[T1]{fontenc}
\usepackage{geometry}
\usepackage{graphicx}
\usepackage{float}
\usepackage{hyperref}
\usepackage[english,indonesian]{babel}

\title{Optimasi Pemrosesan Template LaTeX Menggunakan Pendekatan Berbasis AI}
\author{Academic Support Agent}
\date{\today}

\begin{document}
\maketitle

\begin{abstract}

\end{abstract}

\chapter{PENDAHULUAN}


\chapter{TINJAUAN PUSTAKA}


% Original content (commented out):
% \documentclass{erciyesfk}
\usepackage{etex}
\usepackage[lined,algonl,boxed,commentsnumbered]{algorithm2e}%30.1.14 FEHİM KÖYLÜ
\usepackage[usenames]{color}
\usepackage{amsmath}
\usepackage{amssymb} %30.1.14 FEHİM KÖYLÜ
\usepackage{graphicx,color,psfrag}
\usepackage{enumerate}
\usepackage{rotating}
\usepackage{cite}
\usepackage{psfrag}
\usepackage{pstricks, pst-plot, pst-grad, pstricks-add}
\usepackage{tikz}
\usepackage{pgfplots}
\usepackage{url}
\usepackage{faktor}
\usepackage{lscape}
\usepackage{xcolor}
\usepackage{tabularx}
\usepackage{multirow}
\usepackage{ctable}
\usepackage{float}
\usepackage{afterpage}%30.1.14 FEHİM KÖYLÜ
\usepackage{lineno} %30.1.14 FEHİM KÖYLÜ
\usepackage{subfigure} %30.1.14 FEHİM KÖYLÜ
\usepackage{epstopdf}%30.1.14 FEHİM KÖYLÜ
\usepackage[utf8x]{inputenc} 
\usepackage{algorithmic}%30.1.14 FEHİM KÖYLÜ
\usepackage{setspace}%30.1.14 FEHİM KÖYLÜ
\usepackage{longtable}%02.08.15 FEHİM KÖYLÜ
\usetikzlibrary{positioning,shapes,shadows,arrows}%30.1.14 FEHİM KÖYLÜ
\usepackage{calrsfs}  % 4.6.2015 Kaligrafi paketi F.K.
\DeclareMathAlphabet{\pazocal}{OMS}{zplm}{m}{n} % 4.6.2015 Kaligrafi paketi F.K.
\newcommand{\La}{\mathcal{L}} % 4.6.2015 Kaligrafi paketi F.K.
\newcommand{\Lb}{\pazocal{L}}% 4.6.2015 Kaligrafi paketi F.K. 
\usepackage{array}% 4.6.2015  F.K. 
\usepackage{tablefootnote}  % 27.6.2015 alt yazı paketi F.K.
\usepackage{caption} %18.8 F.K.
\usepackage{amsmath}  %18.8 F.K
\usepackage{listings}
\usepackage{color}

\numberwithin{algocf}{section}  %18.8 F.K.



\ToCLineisDotted %% İçindekiler kısmında sayfa numaralarına kadar olan kısım noktalı olsun
\ToCChapterBold  %% İçindekiler kısmında BÖLÜM 1-2... yazıları bold olsun.
%\TabloListesiYok %% Tablo Listesi kullanılmayacak
%\SekilListesiYok %% Şekil Listesi kullanılmayacak

\setlength{\fboxsep}{0pt}

\input{defination.tex}

\renewcommand{\textfraction}{0.10}
\renewcommand{\topfraction}{0.95}
\renewcommand{\bottomfraction}{0.95}
\renewcommand{\floatpagefraction}{0.35}

\setcounter{totalnumber}{5}

\makeatletter  %18.7.15 FEHİM KÖYLÜ
\renewcommand{\algorithmcfname}{Algoritma}% Update algorithm name
\makeatother




 
\begin{document}	
\input{chapter1.tex} %Özet
\input{chapter2.tex} %Motivasyon
\input{chapter3.tex} %Projenin İçerdiği Yenilik Unsuru
\input{chapter4.tex} %Yöntem ve Metotlar
\input{chapter5.tex} %Proje Planı
\input{chapter6.tex} %Sonuç


\bibliographystyle{erciyesfbefk} 
\bibliography{referanslar}	 
%\appendix
%\input{ek2.tex}

\end{document}


\end{document}